\section{Related Theory}
\label{sec:related}

\sys sits at the intersection of four research lines. For each, we summarize the theory and --- critically for a systems paper --- what each line measures and what it does not.

\subsection{Video analytics systems (the cost-aware lineage)}
\label{sec:related-vas}

The classical lineage optimizes declarative queries over video with cheap-model cascades and indexes. NoScope cascades frame-difference detectors and per-stream proxy CNNs before an expensive target model~\cite{noscope}; Focus builds a coarse ingest-time index and spends the expensive model only on promising index clusters at query time~\cite{focus}; BlazeIt accelerates aggregation and limit queries via control-flow sampling~\cite{blazeit}; TASTI builds a task-agnostic \emph{semantic index} --- embedding buckets whose cluster representatives stand in for full inference~\cite{tasti}; SUPG adds proxy scores with statistical guarantees~\cite{supg}. Resource allocation across streams appears in VideoStorm (quality/resource knob profiles over thousands of streams)~\cite{videostorm}, Chameleon (periodic re-profiling exploiting temporal and cross-camera correlation)~\cite{chameleon}, and VideoEdge (hierarchical cluster placement for camera fleets)~\cite{videoedge}. MIRIS integrates query processing into object tracking~\cite{miris}; Panorama supports unbounded-vocabulary queries over video~\cite{panorama}; VStore manages fidelity tiers of the stored video itself~\cite{vstore}.

This lineage established cost-aware query processing, and we adopt its evaluation discipline (accuracy at a cost target). But it predates open-vocabulary semantics: its queries are fixed-schema (object classes, tracks, counts), its indexes are appearance embeddings, and none of it persists a \emph{memory} that natural-language queries can address. Our TASTI-lite baseline (Section~\ref{sec:exp-tasti}) instantiates this lineage's strongest idea --- the embedding index --- against our text-and-tag memory, and loses on event semantics, consistent with the appearance-signature failure we measured previously (Appendix~\ref{app:e0}).

\subsection{Streaming video LLMs and KV-state memory}
\label{sec:related-kv}

A second line adapts video-LLMs to online streams. VideoLLM-online introduced streaming dialogue training for continuous video~\cite{videollmonline}; Flash-VStream added a memory-based streaming architecture and the RVS benchmarks~\cite{flashvstream}; Dispider disentangles perception, decision, and reaction timing~\cite{dispider}; StreamingVLM sustains infinite commentary with a fixed KV window (512 text tokens + 16\,s of vision) and recency eviction~\cite{streamingvlm}. The KV-retrieval/compression sub-line manages the cache explicitly: ReKV offloads full-fidelity per-frame KV to RAM/disk and retrieves top-$r$ frames per query~\cite{rekv}; LiveVLM compresses online via vision-sink bucketing and retrieves position-agnostic KV pages~\cite{livevlm}; StreamMem prunes to a query-agnostic saliency budget~\cite{streammem}; MuKV compresses at multiple temporal granularities~\cite{mukv}; WeaveTime retrieves past KV via uncertainty triggers~\cite{weavetime}; V-Rex accelerates streaming KV retrieval in hardware~\cite{vrex}. Venus separates memory across an edge--cloud split~\cite{venus}; on the serving side, PagedAttention/vLLM established paged KV management for LLM inference generally~\cite{vllm}, and Stream2LLM addresses scheduling of streaming LLM workloads~\cite{stream2llm}; CodecSight exploits codec signals to cut streaming VLM inference~\cite{codecsight}.

Two properties of this line matter for us. First, its memory is \emph{model-internal state}: not queryable across streams, not portable, and --- at 18.8\,GB/stream-hour for full fidelity at 7B scale~\cite{rekv} --- economically unbounded; ReKV's own appendix names surveillance-scale growth ``unsustainable.'' Second, its evaluation is single-stream and budget-free: token caps are accuracy hyperparameters, and no work in this line reports cost per stream-hour, a fleet model, or an accuracy-vs-cost frontier. We reuse this line as per-stream \emph{components} and as baselines (Sections~\ref{sec:exp-baselines}, \ref{sec:exp-generality}), and our design leaves a slot for an optional KV tier for fine-grained detail (Section~\ref{sec:discussion}).

\subsection{Semantic memory and retrieval over video}
\label{sec:related-semantic}

A third line builds queryable semantic artifacts. VideoRAG constructs a knowledge graph plus text and per-clip visual embeddings over a video corpus, with every 30-s clip captioned and every video transcribed~\cite{videoraghkuds}; an independent NeurIPS variant augments retrieval with OCR/ASR/detection side channels~\cite{videoragneurips}. WorldMM maintains episodic, semantic, and visual memories at four temporal scales, including retaining every sampled frame, driven by frontier API models~\cite{worldmm}. Avas builds an event knowledge graph over continuous streams with agentic tree-search querying~\cite{avas}. The VideoDB report articulates persistent visual memory as layered, analyzer-defined indexes with live streams as first-class sources~\cite{videodb}. LazyVLM decomposes multi-frame queries neuro-symbolically~\cite{lazyvlm}; D\'ej\`a Vu reuses ViT embeddings across temporally redundant frames at query time~\cite{dejavu}; Zelda performs VLM top-$k$ video queries~\cite{zelda}; Lava revisits the classic cascade lineage with language-driven traffic analytics~\cite{lava}; agentic-memory variants (M3-Agent, StreamMeCo, MAGIC-Video) compress long-term agent memory for streaming video~\cite{m3agent,streammeco,magicvideo}.

Every one of these systems writes memory \emph{time-driven and unconditionally}, grows it \emph{unboundedly}, and evaluates \emph{accuracy-only} on a single stream or offline corpus. Avas's throughput (2.5\,FPS indexing vs.\ 2\,FPS input on one RTX\,3090) implies one stream nearly saturates one GPU; VideoDB defers cost measurement explicitly. To our knowledge no published system budgets memory writes, evicts under a storage budget, schedules writes across a fleet, or reports an accuracy-vs-cost frontier. Section~\ref{sec:gap} tabulates this.

\subsection{Video anomaly detection}
\label{sec:related-vad}

Weakly-supervised VAD provides our tier-1 signal and evaluation substrates. UCF-Crime introduced the 13-category surveillance benchmark with 7.8\% anomalous clips~\cite{ucfcrime}; ShanghaiTech offers a 12--13-camera campus setting~\cite{shanghaitech}; XD-Violence adds audio-visual violence detection with 145{,}612 test clips~\cite{xdviolence}. VadCLIP adapts CLIP to weakly-supervised VAD~\cite{vadclip}; cascade designs gate an expensive VLM stage per stream (MemoVAD~\cite{memovad}); Holmes-VAU and VERA bring VLM-based anomaly understanding~\cite{holmesvau,vera}; UCVL benchmarks MLLMs on crime-surveillance tasks~\cite{ucvl}. Our prior measurement study (Appendix~\ref{app:e0}) showed VLM escalation adds $\Delta$AUC\,$\approx 0$ over a strong tier-1 scorer --- which is why, in \sys, the VLM \emph{narrates} gated events rather than \emph{scoring} them, and why category labels come from tier-1 while narrations carry description.

\subsection{Benchmarks for streaming video understanding}
\label{sec:related-bench}

OVO-Bench and its successor OVO-S-Bench evaluate real-time and backward-looking streaming QA on egocentric and web video~\cite{ovobench,ovosbench}; StreamingBench covers 900 videos with 4.5K QA pairs~\cite{streamingbench}; Video-MME, LVBench, MLVU, and LongVideoBench cover offline long-video QA~\cite{videomme,lvbench,mlvu,longvideobench}. None targets always-on \emph{surveillance} streams (no audio, sparse events, multi-camera topology), none carries ground-truth evidence spans suited to memory evaluation, and none is cost-instrumented. This gap motivates \smb (Section~\ref{sec:exp-smb}).
